\subsection{HTTP/1.1 chunked transfer encoding}
\label{sec:bg-chunked}

RFC 9112 \S 7.1 \cite{rfc9112} specifies chunked transfer encoding
as a sequence of length-prefixed chunks. Each chunk has the form
\texttt{<hex-size>;<extensions>$\backslash$r$\backslash$n<data>$\backslash$r$\backslash$n}
where \texttt{<extensions>} is optional and a chunk of size
\texttt{0} terminates the message. A trailer section MAY follow the
zero chunk before the final \texttt{$\backslash$r$\backslash$n}.

The encoding was introduced in RFC 2616 (1999) to support sending
content whose length is not known when the request line is sent.
It is used by streaming clients (cURL with
\texttt{--data-binary @-}, every modern browser forwarding a
streamed body, every LLM SDK streaming a prompt) and by the upload
paths of essentially every public-facing web framework.

The spec defines the \emph{wire} format but does \emph{not}
constrain how a recipient should schedule decoded body bytes into
application code. A recipient may legally:
\begin{itemize}
  \item Deliver each decoded chunk's bytes as a separate
    application-receive event (\emph{per-chunk}; what most servers
    in our sample do).
  \item Coalesce all chunk bytes available in the current TCP
    \texttt{recv()} buffer into one application-receive event
    (\emph{per-recv}; what \fname{System.IO.Pipelines}/Kestrel does).
  \item Fully buffer the entire body up to a configured cap before
    invoking the handler (\emph{fully-buffered}; what Puma's
    Tempfile path does at the handler boundary, but the
    \emph{parser} still loops per chunk).
\end{itemize}

This choice is the load-bearing decision behind everything in this
paper.

\subsection{HTTP server concurrency models}
\label{sec:bg-concurrency}

The servers we measure span the major concurrency paradigms:

\begin{description}
  \item[Event loop, single-thread.] One OS thread, an event loop
    (\texttt{epoll}/\texttt{kqueue}/\texttt{io\_uring}), and
    cooperative tasks. Examples in our sample:
    \fwid{uvicorn} on Python's \fwid{asyncio}, \fwid{hyper}/\fwid{axum}/\fwid{actix-web}
    on Tokio, \fwid{node http.Server} on libuv, \fwid{Vert.x} on Netty,
    \fwid{nginx}, \fwid{HAProxy}.
  \item[Thread per connection.] A worker thread blocks on each
    socket's body read. Examples: Apache \fwid{httpd} \texttt{mpm\_event}
    (per-conn worker), Tomcat (Servlet model), \fwid{gunicorn} sync
    workers, \fwid{waitress}, \fwid{unicorn}.
  \item[N:M scheduler.] Many lightweight tasks (goroutines, fibers,
    BEAM processes) are multiplexed onto fewer OS threads by the
    runtime. Examples: Go \fwid{net/http} (goroutines), \fwid{Cowboy}
    (BEAM processes), \fwid{Falcon} on Ruby (\fwid{async} fibers).
  \item[Actor model.] Message-passing isolation per request.
    Example: \fwid{actix-web} (Actix actor system on top of Tokio).
\end{description}

The per-chunk amplification class does \emph{not} discriminate among
these paradigms. We find it in all four, with broadly comparable
cost constants once normalized for runtime overhead.

\subsection{TCP coalescing and the measurement two-mode dichotomy}
\label{sec:bg-tcp}

A naive measurement of "how slow is the chunked path" depends
critically on whether wire chunks arrive at the server kernel
\emph{one per packet} or \emph{batched many per packet}. Both
scenarios are realistic:

\begin{itemize}
  \item In a kubernetes pod-to-pod call across a docker bridge, the
    kernel will coalesce dozens to hundreds of small TCP segments
    into one \texttt{recv()} of $\sim$\SI{4}{\kibi\byte}.
  \item In a slow-drip / slow-loris-style attack, or behind a CDN
    that does not aggregate, each chunk can land in its own
    \texttt{recv()}.
\end{itemize}

The two regimes can differ by orders of magnitude on the same
server. We therefore measure both, denoting them \emph{Mode A}
(bridge-coalesced) and \emph{Mode B} (paced, each chunk in its
own \texttt{recv()}); see Section~\ref{sec:methodology}.

\subsection{Slow-loris lineage and CWE-407}
\label{sec:bg-cwe}

The attack pattern of "send small request fragments slowly to keep
a server busy" goes back to Apache Killer / slow-loris and the RUDY
family. The class is catalogued as CWE-407 (Inefficient Algorithmic
Complexity)~\cite{cwe407}; recent ecosystem-relevant CVEs include
CVE-2025-67725 (Tornado quadratic header
coalescing)~\cite{cve_2025_67725_tornado}, CVE-2025-14550
(Django ASGI repeated-header concatenation)~\cite{cve_2025_14550_django},
CVE-2026-7790 (cowlib chunk-size hex
amplification)~\cite{cve_2026_7790_cowlib}, CVE-2023-39326 (Go
net/http chunk-extension cost cap)~\cite{cve_2023_39326_go}, and
the older CVE-2014-0075 (Tomcat chunk-size integer
overflow)~\cite{cve_2014_0075_tomcat}.

Each of these CVEs targets a \emph{specific} cost in the
chunked-TE pipeline (per-byte hex parsing, per-header coalescing,
integer-overflow on chunk size). None targets the per-chunk
application-dispatch cost we measure here. The bug class we
characterize is the same family but a previously-unnamed shape;
Section~\ref{sec:related} positions this lineage against the
broader HTTP DoS literature (slow-DoS, algorithmic-complexity, and
the modern HTTP/2 control-plane and decompression attacks).
